\section{統計の限界・バイアス・倫理 — 手法の外側を知る}

\subsection{1章で告げた「正しく計算しても間違う」の正体}

第2章から第8章まで, データを要約し, 形を見て, 幅を測り, 差を吟味し, 関係を測り, モデルにしてきた.
道具は揃った. ここまでくれば, 数字を出して結論を言う作業はだいたいできるはずだ.

ところが, ここまで学んだ道具を全部正しく使っても, 結論が間違うことがある.
1章のはじめで, この種の例が世の中には存在する, とだけ伝えてあった.
あの予告の答え合わせを, ここでやる.

1973年, カリフォルニア大学バークレー校の大学院は, 性別を理由に女子受験生を差別しているとして訴えられた.
証拠とされたのは合格率の差だ.
全志願者を集計すると, 男子の合格率は約44\%, 女子は約35\%.
9ポイントの差は無視できない.

この数字だけ見れば, 大学院に問題があると考えるのが自然だ.
ところが大学院側が学部別に集計し直したところ, 6つの主要学部のうち4つで, 女子の合格率は男子と同じか, むしろ高かった.
学部単位ではどこにも女子不利の傾向はない. それなのに大学全体の合計だけが男子有利に出る.

なぜそうなったのか. 応募の分布が学部によってまるで違ったのだ.
女子は合格率がそもそも低い学部（人文系で20\%から30\%）に多く志望し, 男子は合格率が高い学部（理工系で60\%から70\%）に多く志望していた.
各学部の中では公平でも, 女子が「狭き門の学部」に集中していたために, 全体平均では女子の合格率が下に引かれた.

集計の単位を変えただけで結論が逆転する.
このような現象を\textbf{シンプソンのパラドックス}と呼ぶ.
平均は2章で, 散布や相関は7章で扱った. どの計算も正しい.
なのに, 全体集計と層別集計で正反対のことが言える.

ここで読者に頼みたいのは, この種の食い違いを「珍しい例外」と片づけないことだ.
第3変数が群によって違う割合で混ざっていれば, 同じ仕掛けはどんな分析でも起こりうる.
2章から8章で身につけた計算は, 「集計の単位は妥当か」「混ざっているはずの第3変数はないか」という問いとセットでなければ, 簡単に道を踏み外す.

8章の最後で, 「手法の外側を知ることが, 手法を正しく使うための最後の一歩になる」と書いた.
手法の外側とは, おおまかに言って3つの方角だ.
標本をどう集めたか. 集まったデータの中で何が見えていないか. 出てきた数字をどう解釈したか.
本章ではこの3方向の落とし穴を順に整理し, 続いて手法そのものが暗黙に置いている前提に踏み込む.
最後に, 数字を扱う側の責任の話で本書を閉じる.

まずは標本の集め方から始めよう.
3つの落とし穴のうち, もっとも身近で, もっとも誤解されやすいのがこの入口だからだ.


\subsection{バイアスの三つの顔 — 標本・観測・解釈}

データに歪みが入る場面は, 大きく分けて3つある.
標本を集める段階, 集めたデータの中身が偏っている段階, 集まったデータを解釈する段階.
データの流れの上流・中流・下流に, それぞれ違う種類の歪みが棲んでいる.

\subsubsection*{上流 : サンプリングバイアス}

5章で信頼区間を計算したとき, 「件数を増やせば幅は狭くなる」という結果を手に入れた.
$n$を増やせば$\sqrt{n}$で精度が上がる. 計算式の上ではその通りだ.
ただし, この話は標本が母集団から無作為に抽出されているという前提の下でだけ成り立つ.

抽出が偏っていたらどうなるか.
たとえば, 関心のある母集団は「日本に住む有権者」なのに, 集めたのは「平日の昼に渋谷駅にいた人」だとしよう.
件数を10万人に増やしたところで, 平日の昼に渋谷駅にいない人の意見はゼロのままだ.
標本平均が近づいていく先は, 「日本に住む有権者」ではなく「平日の昼に渋谷駅にいた人たち」の真の平均でしかない.
4章のCLTで「標本平均は母平均に近づく」と言ったとき, この母平均は標本を取った母集団の平均であり, 知りたかった母集団のそれとは限らない.

これが\textbf{サンプリングバイアス}である.
標本が母集団を代表していないとき, $n$を増やすほど狭く, しかも誤ったところに信頼区間が落ちていく.
件数で殴って解決できる問題ではない.

5章で紹介したリテラリー・ダイジェスト誌の話を思い出してほしい.
1936年の大統領選で, 同誌は約240万人を対象とした大規模調査からランドンの勝利を予測した.
ところが結果はルーズベルトの圧勝だった. 同じ年, ジョージ・ギャラップは5万人の調査で正しく予測している.
件数では数十倍劣るギャラップが勝ったのは, 標本の代表性で勝っていたからだ.
ダイジェスト誌のリストは電話帳と自動車登録名簿で, 大恐慌下で電話も車も持たなかった貧困層をまるごと外していた.
あの逸話は, サンプリングバイアスの教科書的な実例にあたる.

\subsubsection*{中流 : 生存者バイアス}

標本の代表性は, 「取れたデータの中身」を眺めても判断できないことがある.
そもそも取れなかったデータがあると気づかないと, 結論は静かに歪む.

第二次大戦中, 米軍は帰還した爆撃機の被弾箇所を調べて, どこを補強すべきかを検討していた.
集計してみると, 翼の先端と機体後部に弾痕が集中している.
最初の判断は素直だった. 「弾痕が多い場所を補強しよう」.

統計学者アブラハム・ワルドはこの結論を止めた.
「補強すべきは, 弾痕の少ない場所のほうだ」.
帰還した機は, 翼や機体後部に被弾しても飛び続けられた機だ.
エンジンや操縦席にやられた機は, そもそも基地に戻ってこない.
分析の対象に上がっていたのは「生き残った機」だけで, 「落ちた機」の被弾箇所はデータに現れていなかった.

これが\textbf{生存者バイアス}である.
「観測できたものだけ」を見て結論を出すと, 観測できなかったものが体系的に抜け落ちる.
事業継続中の企業だけを集めて成功要因を語る研究, 在籍中の社員だけを集めた満足度調査, 完走したマラソン選手だけを集めた走法分析.
どれも, 退場した側に重要な情報が眠っている可能性がある.

データの中身を眺めるだけでは, 何が抜けているかは見えない.
「このデータには本当はどんな対象が含まれるべきか, そのうち実際に観測できたのは誰か」を, データを見るのとは別の頭で考える必要がある.

\subsubsection*{下流 : 確証バイアス}

仮にデータの集め方が完璧で, 観測の漏れもなかったとしよう.
それでも, 出てきた数字をどう解釈するかの段階で歪みが入る.

人間は, 自分の仮説に合うデータを重く受け取り, 合わないデータを軽く流す癖を持っている.
これを\textbf{確証バイアス}と呼ぶ.
心理的な現象として聞こえるが, 統計の道具で正当化された確証バイアスもある.

公平なコインを20枚用意して, それぞれ100回投げ, 各コインで「表と裏の出方が偏っているか」を有意水準5\%で検定すると考えてみよう.
20枚すべてが本当に公平だったとしても, 偶然どれか1枚は「有意」と出る.
具体的には, どこかで偽陽性が出る確率は$1 - (1 - 0.05)^{20} \approx 0.64$, つまり約64\%だ.
20回試せば, 6割以上の確率で「有意な発見」がでっち上がる.

ここで実験者が「有意になった1枚」だけを論文にし, 残り19枚を引き出しにしまったらどうなるか.
表に出るのは「特殊なコインを発見した」という主張だけだ.
6章で習ったp値の計算は, どの1枚についても正しく行われている.
それでも結論は嘘になる.

これが\textbf{p-ハッキング}と呼ばれる行為である.
6章で覚えたp値が, ここでは正当化の道具として悪用されている.
都合のいい結果が出るまで分析の角度や対象を変え続けると, 確証バイアスは自動化される.

ここまで3つの歪みを並べた.
標本の集め方（上流）, 観測対象の選別（中流）, 解釈の仕方（下流）.
データが流れる順番に沿って, 別々の場所に違う罠が仕掛けられている.

3類型は別々の知識ではなく, 1本の流れの上の地点だ.
分析を始める前にまず上流を疑い, 手元のデータを見ながら中流を疑い, 結果を読むときに下流を疑う.
本章の残りはこの順序のまま進む.
次は, 集まったデータに手法を当てる段階, つまり「中流から下流の入口」に位置する話に踏み込む.


\subsection{前提が崩れるとき — 独立性・線形性・等分散性}

データの集め方や選別の歪みを乗り越えたとしよう.
それでも, 手法そのものが暗黙に置いている前提が崩れていれば, 出てきた数字の意味は変わる.
2章から8章で使ってきた手法のうち, 推論に関わる部分は概ね3つの前提に乗っている.
データ点が独立であること. 関係がおおむね直線で表せること. 残差の散らばりがどの場所でも同じくらいであること.

\subsubsection*{独立性 : 隣の点と関係していないか}

4章のCLTの前提として「独立同分布」を導入した.
ここでの「独立」は, データ点どうしが互いに無関係であることを指す.
ある月のデータが翌月のデータに影響していたら, 独立ではない.

8章で扱った広告費と売上の例を思い出してほしい.
あの10ヶ月分のデータは時系列に並んでいた.
回帰の計算式は, 各データ点を独立に扱って残差の二乗和を最小化している.
ところが, 先月の広告効果が今月にも残っていたとしたらどうなるか.
ある月の売上が高ければ, その勢いで翌月も高い. 残差の符号は隣どうしで似てしまう.

このとき計算式は止まらない. 傾きも切片も決定係数も, 数字としては出る.
ただし, 出てきたSEは「独立だった場合に成り立つ式」で計算されているから, 独立でない実態に対しては不当に小さい値になる.
SEが小さいと, 信頼区間も狭く出る. 6章のp値も小さく出る.
本当は不確かな結論が, 紙の上では確実そうに見える.

時系列データだけではない. クラスごとの生徒, 同じ会社の社員, 同じ地域の住民. 集団内の点は互いに似ている傾向がある.
集団効果を無視して全員を独立だと扱うと, 同じ罠が起こる.
独立性は, 計算が動くかどうかではなく, 計算結果の信頼性を決めている.

\subsubsection*{線形性 : 直線で表せる関係なのか}

8章で残差プロットを導入したとき, アンスコムの第2組の話をした.
4組のデータが$r \approx 0.82$, $R^2 \approx 0.67$という似た数字を返したが, 第2組は実際にはきれいな放物線で, 直線を当てはめること自体が無意味だった.

直線モデルは, 「関係がおおむね直線で表せる」という前提の下で意味を持つ.
本当の関係が放物線, 指数, 周期的なものだったら, 直線を引いても得られるのは「もっとも近い直線」だけだ.
傾きや切片の数値は出るが, それらは現実の関係を要約していない.

線形性が崩れているかどうかは, 残差プロットで目視できる.
予測値の小さい所と大きい所の両方で残差が正に偏り, 中央で負に偏っていれば, 関係はU字型だ.
逆向きのパターンも同様の警告だ.
8章で「ゼロのまわりにランダムに散らばっていればよい」と書いたのは, この前提が満たされているかを目で確かめる手順だったわけだ.

線形性の確認は, 数値だけ見ていても気づけない.
$R^2$が0.7を超えていても, 残差プロットが系統的なパターンを示していれば, 直線モデルは不適切だ.
逆に$R^2$が0.4でも, 残差にパターンがなければ, それはそれで意味のある直線関係である.
$R^2$は当てはまりの「総量」しか語らず, 「外し方の質」は残差プロットでしかわからない.

\subsubsection*{等分散性 : 場所によって散らばりが違わないか}

3つめの前提が, 残差の散らばりがデータの位置によって変わらないこと, つまり等分散性である.
6章のt検定でWelch法を採ったのは, 2群間で分散が違うことを許す扱いを選んだという意味で, この前提を緩めた選択だった.

回帰では事情が少し違う.
広告費が大きい月ほど売上の上下幅も大きくなる, というように, 予測値の大きさによって残差の散らばりが変わる場合がある.
残差プロットでは, 左端は密で右に行くほど縦方向に広がっていく扇形が見える.

このとき, 直線そのものは間違っていない.
平均的な傾きは正しく推定されている.
ただ, 「来月の広告費を50万円にしたら売上はおよそいくつか」と幅つきで予測したい場合, 場所ごとに本当の散らばりが違うのに一律のSEで予測区間を出すと, 区間が小さい場所では狭すぎ, 大きい場所では広すぎる.
平均値の推定は使えても, 予測の幅の保証は揺らぐ.

ここまでの3つの前提を点検する手段は, ほとんど残差プロット1枚に集約される.
残差を予測値に対してプロットし, ランダムに散らばっていればよし, 系統パターンや扇形が見えたら警告.
8章で導入したあのプロットは, 単なる確認作業ではなく, 線形回帰という道具を使ってよいかどうかを判断する診断装置だったのだ.

ここまでで, データの集め方と手法の前提は点検する目を持った.
それでもまだ, 「正しく計算しても結論が間違う」場面がある.
計算自体は正しく動き, 前提も満たされている. それなのに, 結果の読み方が間違っているケースだ.
次節はその話に移る.


\subsection{正しく計算しても結論が崩れるとき — p値・相関・外挿}

ここまでは, データの上流から手法の前提までの話だった.
すべての段階を正しく済ませたとしよう.
それでも, 出てきた数字の解釈が間違っていれば, 結論はやはり間違う.
誤読が起きやすいのは, おおまかに3か所だ. p値, 相関と因果の関係, データ範囲の外側.

\subsubsection*{p値の3つの誤読}

6章でp値を「帰無仮説が正しいと仮定したとき, 観測値以上に極端な結果が出る確率」と定義した.
この定義はそれなりに込み入っているせいで, 短く要約しようとすると意味がずれる.
よく聞く誤読を3つ挙げておこう.

ひとつ. 「p値は帰無仮説が正しい確率である」という読み方.
これは違う. p値は「帰無仮説が正しいと仮定したとき」の珍しさを測っているのであって, 帰無仮説そのものの真偽の確率を返してくれるわけではない.
p値が0.03でも, 仮説が間違っている確率は3\%ではない.

ふたつ. 「p<0.05なら効果は大きい」という読み方.
p値は珍しさの指標であって, 効果の大きさは別の指標で測る.
たとえば2群間の平均差を比べる場面で, 群1の平均が3.04, 群2が3.01だったとしよう.
差は0.03点だ. 各群の標本数が10000ある場合, この差はおそらくp<0.05で有意になる.
件数が大きいと, 実用上は無視できる差でも「珍しい」と判定されやすい.
逆に, 件数が少なければ, 大きい差でも有意にならないことがある.

みっつ. 「p>0.05なら効果はない」という読み方.
これも違う. 「珍しいと言える証拠が足りない」と「効果がないことが証明された」は別の話だ.
n=20で実験して有意にならなかったからといって, 効果が存在しないわけではない.

p値の話を扱うときは, p値だけで報告を済ませない方がよい.
5章で覚えた信頼区間と, 6章Part Bで触れた効果量を一緒に出す.
信頼区間は「効果がだいたいどの範囲にありそうか」, 効果量は「差そのものの大きさ」, p値は「偶然で説明しにくいかどうかの目安」.
3つはそれぞれ別の問いに答えている.

報告の側もこの3点をセットで書き, 読み手の側もこの3点を求めるとよい.
p値だけが踊っている結論は, それだけで疑う材料になる.

\subsubsection*{相関は因果ではない, ではどう違うのか}

7章と8章で「相関は因果を意味しない」と何度か言った.
言葉としては理解しやすいが, 「ではどう違うのか」と聞かれると, 案外答えにくい.
相関が因果でない場合の典型を3つ並べておく.

第1は\textbf{交絡}.
アイスクリームの月別消費量と水難事故の月別件数には, 強い正の相関がある.
アイスを食べると溺れるわけではない. 気温が高い月にアイスもよく売れ, 同じく気温が高い月に水場へ行く人が増えて事故も増える.
両方を独立に押し上げている第3変数（気温）があるとき, アイスと水難事故の間に直接の因果がなくても相関は出る.
この第3変数を交絡因子と呼ぶ.

ヨーロッパ各国でコウノトリの巣の数と人間の出生数に正の相関がある, という有名な話もこの仲間だ.
コウノトリが赤ん坊を運んでくるわけではない. 国の面積や人口といった「国の規模」が両方を押し上げているだけだ.
交絡因子は気温のような物理量だけでなく, 国の規模のような構造量でも起こる.

第2は\textbf{逆因果}.
$x$が$y$を動かしているのではなく, $y$が$x$を動かしている可能性がある.
例えば「医療支出が多い地域ほど健康指標が悪い」というデータが出たとしよう.
医療支出が悪さをしているのではなく, 健康状態が悪い地域だからこそ医療支出が増えている, という方向の因果かもしれない.
矢印の向きを取り違えると, 結論はそのまま逆転する.

第3は\textbf{偶然の一致}.
無関係な2つの量を大量に集めて相関を計算すれば, そのうちのいくつかは偶然強く相関する.
「米国のチーズ消費量とベッドシーツでの絞死者数」のような, ふざけた相関集が世間にあるくらいだ.
たくさん試せば偽陽性は当然出る, という前節のp-ハッキングと同じ構造の話である.

相関を見たときに自動的に問うべき3つは決まっている.
交絡しそうな第3変数はないか.
因果の向きはどちらか.
似たような量を大量に試した結果ではないか.
この3つに即答できないなら, 「強い相関がある」というだけの結論で立ち止まらないほうがよい.

\subsubsection*{データの外側に踏み込むとき}

8章で広告費と売上の回帰直線$y \approx 119 + 2.3x$を引いた.
「来月の広告費を50万円にしたら売上はおよそ234万円」と予測した.
50万円はデータの範囲（20万円から80万円）に収まっているから, この予測には根拠がある.

では, 広告費200万円ならどうか.
式に代入すれば579万円だ. 計算は問題なく動く.
しかし200万円はデータの範囲をはるかに超えた値である.
広告費20万円から80万円の範囲で観測された直線関係が, 200万円まで続いている保証はどこにもない.
広告には効果の頭打ちがある場面が多いし, 別のメカニズムが働き始める場合もある.

これが\textbf{外挿}の問題だ.
データの範囲外に直線を延ばして予測すると, 「直線が続いている」という根拠のない仮定を黙って置くことになる.
電卓は何も警告しない. 警告するのは使い手だ.

切片の解釈も外挿と地続きの話だ.
$a \approx 119$は, 式の上では「広告費がゼロのときの売上119万円」となる.
しかし広告費ゼロの月は実データにない.
切片の119万円は, 直線の位置を決めるために計算上現れた数値であって, 「広告ゼロでも119万円売れる」ことを意味しない.

「数字が出る」と「現実に当てはめてよい」は別だ.
モデルが描く直線は, 観測した範囲で関係を要約する道具であって, その外側まで現実を保証する装置ではない.

ここまでが, 計算は正しいのに結論が間違う典型だった.
最後は, 数字を扱うこと自体に伴う責任の話に視野を広げて, 本書を閉じる.


\subsection{データを扱う責任と, この本の終わりに}

ここまでは技術と解釈の話だった.
最後に1つ, 技術の続きとしての話をしておく.

倫理という言葉を聞くと, 技術書から少し離れた話に聞こえるかもしれない.
だが本章で並べた歪みは, 個人がうっかり犯すだけで終わらず, 社会の規模で誰かに不利益を与える形にもなりうる.
そこに歯止めをかけるのが, データを扱う側の責任の話だ.

\subsubsection*{設計の逸脱が社会保障に直結した例}

日本では2018年に, 厚生労働省の毎月勤労統計をめぐる問題が発覚した.
毎月勤労統計は雇用保険や労災保険の給付額を決める基礎資料の1つであり, 法律上は従業員500人以上の大規模事業所は全数調査することになっていた.
ところが東京都分について, 2004年から長年, 全数調査ではなく一部の抽出調査が行われていた.
さらに, 抽出に伴って必要な復元処理（拡大処理）が抜けていた.

結果として平均賃金が体系的に低めに計上され, 雇用保険などの給付額も低めに支給された.
影響を受けた人は延べ約2{,}000万人, 過少支給の総額は約500億円超と公表されている.
追加給付のための事務費は別途数百億円規模で発生した.

「全数調査するはずの場所で抽出をしていた」という設計の逸脱が, 社会保障の配分という実体に直結した事例だ.
9.2で扱ったサンプリングバイアスが, 個人の研究の話ではなく公的な仕組みの中で起きていたわけだ.
倫理の話は, 「他人事の高尚な話」ではなく, 「この本で学んだ歪みが社会の規模で起きないよう仕組みで抑える」という地続きの話なのだ.

\subsubsection*{最低限の3点 : 同意と目的, 保護, 透明性}

データを扱うときの規範は, 国や分野によって細部が違うが, おおむね次の3点に集約できる.

ひとつめは, 集めるときに目的を明示し, 同意を得ること.
何のためのデータなのか, どこまで使われるのかを, データを提供する側が理解できる形で伝える.
あとから目的外に流用しない.

ふたつめは, 個人が特定できる情報を保護すること.
氏名や住所そのものを直接保管するかどうかの話だけではない.
複数の項目を組み合わせると個人が特定できてしまう場合があるので, 公開する集計の単位や粒度には注意がいる.

みっつめは, 報告の透明性. 手順, 前処理, 除外したデータの基準を, 第三者が検証できる形で公開する.
9.2で扱った確証バイアスやp-ハッキングは, 透明性の欠如の上にこそ繁茂する.
誰がどの段階で何を選んだかが見えていれば, 外側からの検証によって歪みは正されうる.

3点はバラバラのルールではない.
9.2で見た上流・中流・下流の歪みを, 社会の側で抑える装置だと考えればよい.
集め方を歪ませない（上流対策）, 個人を巻き添えにしない（中流対策）, 解釈の自由を残しすぎない（下流対策）.
本章で習った歪みの地図と, ここでの3点は, 同じ地形を別の角度から眺めた地図である.

\subsubsection*{ここまでに何ができるようになったか}

本書を1章から振り返ると, ずいぶん遠くまで来た.

2章でデータの大きさと散らばりを言葉にした. 平均と中央値, 分散と標準偏差, 四分位範囲.
3章でデータの形を見る道具を手にした. ヒストグラム, 箱ひげ図, 散布図.
4章で「取り直したらばらつく」という事情を引き受け, CLTを通じて正規分布が現れる理屈を理解し, SE = $s/\sqrt{n}$で推定値の幅を測れるようになった.
5章では, その幅をどう読み, どう使うかに集中した. 95\%信頼区間の正しい読み方, 平均・割合・差の3ケース, 件数を増やすと幅は$\sqrt{n}$倍ずつ縮むこと.
6章で「偶然では説明しにくいかどうか」を数で確かめる枠組みを得た. 帰無仮説と対立仮説, 検定統計量, p値, 有意水準. Part Bでt検定とカイ二乗検定の数値計算と効果量に踏み込んだ.
7章で2変数の関係の方向と強さを相関係数で測った. 8章では関係を直線で表し, $R^2$と残差プロットで自分のモデルを疑う習慣をつくった.
そして本章で, 手法の外側にある歪みを地図にした.

ここまでに獲得した能力を一言で表すなら, 「数字が出る」から「数字を疑える」までの距離だ.
学び始めの頃は, 計算結果を出せるようになることが目標だったかもしれない.
ここまで来た今, それは出発点にすぎなかったとわかる.
出した数字の前提は何か. その前提は満たされているか. 集計の単位は妥当か. 第3変数を見落としていないか. データの範囲外まで延ばして語っていないか.
数字を出した後にこれだけ問えるようになっていれば, 本書の役目はおおむね果たされた.

\subsubsection*{次の一歩}

本書は統計の骨格に絞った. 補足概念として各章の末尾に名前だけ挙げた話題は, それぞれ独立した教科書1冊分の世界が広がっている.
多重比較の制御, ブートストラップによる区間推定, 重回帰と多重共線性（VIFの議論）, 時系列分析の独立性の扱い方, 事前登録と再現性の手順, データガバナンスとプライバシ.
どれも, 本書で身につけた骨格の上に積み上げる種類の話だ.
興味の向く順に, 必要を感じたところから手を伸ばせばよい.

最後に1つだけ提案する.
読者自身の手元のデータで, 本書の手順を最初から最後まで一度回してみてほしい.
要約し, 形を見て, 幅を測り, 差を吟味し, 関係を見て, モデルにして, 出した結論を疑う.
その過程で, 本書のどこかが「もっと知りたい」「ここはなぜそうなのか」と引っかかるはずだ.
その引っかかりこそが, 統計学の次の章への入口になる.

教科書の終わりは, 学びの終わりではなく, 自分の問いに戻る場所だ.
ここから先は, あなたが自分のデータと向き合う番である.
